\documentclass[10pt,a4paper]{amsart}
\usepackage[margin=19mm]{geometry}
\usepackage{amsmath,amssymb,mathtools}
\usepackage[scr=rsfs,cal=euler]{mathalfa}
\usepackage{xfrac}
\usepackage[unicode,hidelinks,bookmarks=false]{hyperref}
\newcommand{\ie}{i.e.\ }
\newcommand{\cf}{cf.\ }
\newcommand{\resp}{respectively\ }
\newcommand{\Subsection}{\subsection}
\providecommand{\nobreakdash}{\nobreak\hbox{-}}
\setlength{\emergencystretch}{2em}
\multlinegap0pt
\usepackage{amssymb}
\usepackage{mathtools}
\usepackage{xcolor}
\newtheorem{lemma}{Lemma}
\newcommand{\R}{\mathbb{R}}
\newcommand{\dd}{\,\mathrm{d}}
\newcommand{\ellf}{\ell}
\title{Persistence and analyticity of plasmon poles for the linearized Coulomb--Hartree--Fock equation in the small-exchange regime}
\author{Yuya Dan}
\date{Source: arXiv:2608.29671v1 (CC BY 4.0)}
\begin{document}
\maketitle

\noindent\emph{Source and licence.} Persistence and analyticity of plasmon poles for the linearized Coulomb--Hartree--Fock equation in the small-exchange regime. Version arXiv:2608.29671v1 (CC BY 4.0), distributed by arXiv under the Creative Commons Attribution 4.0 International licence, http://creativecommons.org/licenses/by/4.0/, as recorded in the arXiv licence metadata for this exact version and shown on the abstract page https://arxiv.org/abs/2608.29671v1. Full licence notice: \texttt{licenses/arxiv-2608.29671-CC-BY-4.0.txt}.\par\smallskip

\noindent\textit{Editorial scope.} Counted unit: the article's lemma lem:AXident (source0.tex lines 770--778). The article's complete original proof of it is delivered with it (source0.tex lines 780--791, one \texttt{proof} environment of 12 lines). No mathematics is written, repaired or re-derived: the statement and the proof are the article's own lines, in the article's own notation, and no retained line is substituted or re-flowed.  Retained uncounted as the source-local context the proof needs: lines 136--139.  Nothing was omitted from any retained span, so the package carries no omission record.  No retained line contains a citation, so the wrapper carries no bibliography and no bibliography entry.  No retained line contains a figure environment, an image-inclusion command or drawing code, so no image is delivered and the package declares no asset dependency.  Editorial formatting only, in addition: an AMS \texttt{amsart} preamble that restates the article's own declarations and macros used by the retained lines, namely the environments \texttt{lemma} on separate counters, together with the standard packages those lines need.  The retained environments print in this document as Lemma 1 in order of appearance, whereas the article prints the same items with the numbering of its own full document.  The complete native source of the article as distributed in the arXiv e-print for this exact version is bundled unchanged as \texttt{sources/208876/source0.tex}.
\begin{equation}\label{eq:selfenergy}
\Sigma(p)=\frac{1}{(2\pi)^3}\int_{\R^3}V(p-q)\,g(q)\dd q,\qquad
\epsilon_\eta(p)=|p|^2-\eta\,\Sigma(p).
\end{equation}

\begin{lemma}[The self-energy increment is the Coulomb potential of $a_k$]\label{lem:AXident}
$A_{X,k}=W a_k$. Consequently, with $L=L_{0,*}$ and $\operatorname{Diff}(r):=\ellf(L^{-1}A_{X,k}L^{-1}a_k)-\ellf(L^{-1}a_kW_RL^{-1}a_k)$,
\begin{equation}\label{eq:diffcombined}
\operatorname{Diff}(r)=\frac{2r^2}{(2\pi)^6}\iint_{\R^3\times\R^3}
V(p-q)\,a_k(p)\,a_k(q)\,\frac{(p_1-q_1)^2}{L(p)^2L(q)^2}\dd p\dd q ,
\qquad k=re_1 ,
\end{equation}
an absolutely convergent integral with bounded integrand \textup{(}the factor $(p_1-q_1)^2$ cancels the Coulomb singularity, since $V(p-q)(p_1-q_1)^2\le4\pi$\textup{)}.
\end{lemma}

\begin{proof}
Shifting the integration variable in \eqref{eq:selfenergy},
\[
A_{X,k}(p)=\Sigma\Bigl(p-\tfrac k2\Bigr)-\Sigma\Bigl(p+\tfrac k2\Bigr)
=\frac{1}{(2\pi)^3}\int V(p-q)\Bigl[g\Bigl(q-\tfrac k2\Bigr)-g\Bigl(q+\tfrac k2\Bigr)\Bigr]\dd q=(Wa_k)(p).
\]
Hence both terms of $\operatorname{Diff}$ are quadratic forms with the same kernel:
\[
\operatorname{Diff}(r)=\frac{1}{(2\pi)^6}\iint V(p-q)\,a_k(p)a_k(q)\Bigl[\frac1{L(p)^2}-\frac1{L(p)L(q)}\Bigr]\dd p\dd q .
\]
Symmetrizing in $(p,q)$ and using $\frac12\bigl[\frac1{L_p^2}+\frac1{L_q^2}-\frac2{L_pL_q}\bigr]=\frac{(L_p-L_q)^2}{2L_p^2L_q^2}$ with $L(p)-L(q)=2r(p_1-q_1)$ gives \eqref{eq:diffcombined}. (On $S_k$, $W$ and $W_R$ coincide when applied to functions supported in $S_k$.)
\end{proof}

\end{document}
